\section{Second-Order Equations and Oscillation}
\label{sec:ch4-second-order}

Second-order equations arise naturally because acceleration is the second derivative of position. The prototype is the oscillator, whose mathematics later reappears in normal modes, waves, quantum mechanics, and field theory.

\subsection{Constant-coefficient equations}
Consider
\[
a\ddot x+b\dot x+cx=0.
\]
Trying $x=e^{rt}$ gives the characteristic equation
\[
ar^2+br+c=0.
\]
Its roots determine the qualitative motion: distinct real roots give sums of exponentials, a repeated root introduces a factor of $t$, and complex-conjugate roots produce oscillation.

\subsection{The ideal harmonic oscillator}
Hooke's law $F=-kx$ and Newton's second law produce
\[
m\ddot x+kx=0.
\]
Defining $\omega_0=\sqrt{k/m}$ gives
\[
\ddot x+\omega_0^2x=0,
\]
with solution
\[
\boxed{x(t)=A\cos(\omega_0t)+B\sin(\omega_0t)}.
\]
Equivalently, $x(t)=C\cos(\omega_0t-\phi)$. The constants encode amplitude and phase.

The velocity is $v=\dot x$, and the energy
\[
E=\frac12m\dot x^2+\frac12kx^2
\]
is conserved because
\[
\frac{dE}{dt}=\dot x(m\ddot x+kx)=0.
\]

\begin{figure}[htbp]
\centering
\begin{tikzpicture}[x=.8cm,y=1.1cm]
\draw[->] (0,0)--(7.2,0) node[right] {$\omega_0 t$};
\draw[->] (0,-1.4)--(0,1.5);
\draw[chapterblue,thick,domain=0:6.6,samples=120] plot (\x,{cos(57.2958*\x)});
\draw[orange!70!black,thick,dashed,domain=0:6.6,samples=120] plot (\x,{-sin(57.2958*\x)});
\node[chapterblue] at (5.8,1.1) {$x/A$};
\node[orange!70!black] at (5.8,-1.05) {$v/(A\omega_0)$};
\end{tikzpicture}

\caption{Position and velocity of an ideal oscillator. Their quarter-cycle phase offset encodes continual exchange between potential and kinetic energy.}
\label{fig:ch4-harmonic-motion}
\end{figure}

\subsection{Damped motion}
With linear resistance $-b\dot x$,
\[
m\ddot x+b\dot x+kx=0.
\]
Define
\[
\gamma=\frac{b}{2m},\qquad \omega_0=\sqrt{\frac{k}{m}}.
\]
The characteristic roots are
\[
r=-\gamma\pm\sqrt{\gamma^2-\omega_0^2}.
\]
This produces three regimes:
\begin{itemize}
\item underdamped: $\gamma<\omega_0$;
\item critically damped: $\gamma=\omega_0$;
\item overdamped: $\gamma>\omega_0$.
\end{itemize}
For underdamping,
\[
x(t)=e^{-\gamma t}\left[A\cos(\omega_dt)+B\sin(\omega_dt)\right],
\qquad \omega_d=\sqrt{\omega_0^2-\gamma^2}.
\]

\begin{figure}[htbp]
\centering
\begin{tikzpicture}[x=.72cm,y=1.0cm]
\draw[->] (0,0)--(8.2,0) node[right] {$t$};
\draw[->] (0,-1.25)--(0,1.35) node[above] {$x$};
\draw[chapterblue,thick,domain=0:8,samples=160] plot (\x,{exp(-0.28*\x)*cos(100*\x)});
\draw[orange!70!black,thick,dashed,domain=0:8,samples=100] plot (\x,{(1+\x)*exp(-\x)});
\draw[green!45!black,thick,dotted,domain=0:8,samples=100] plot (\x,{1.15*exp(-0.35*\x)-.15*exp(-2.2*\x)});
\node[chapterblue,anchor=west] at (5.7,.62) {underdamped};
\node[orange!70!black,anchor=west] at (5.7,.34) {critical};
\node[green!45!black,anchor=west] at (5.7,.08) {overdamped};
\end{tikzpicture}

\caption{Underdamped, critically damped, and overdamped return toward equilibrium.}
\label{fig:ch4-damping-regimes}
\end{figure}

\subsection{Driven oscillation and resonance}
An external sinusoidal force gives
\[
m\ddot x+b\dot x+kx=F_0\cos(\Omega t).
\]
After transients decay, the steady-state response has the form
\[
x_{\mathrm{ss}}(t)=A(\Omega)\cos(\Omega t-\delta),
\]
where
\[
A(\Omega)=\frac{F_0/m}{\sqrt{(\omega_0^2-\Omega^2)^2+(2\gamma\Omega)^2}}.
\]
The response is largest near resonance. Damping limits the peak and produces a frequency-dependent phase lag.

\begin{keyidea}[Homogeneous plus particular]
For a linear driven equation, the complete solution is
\[
x=x_{\mathrm{transient}}+x_{\mathrm{steady}}.
\]
The homogeneous part remembers the initial condition; damping erases that memory. The particular part follows the persistent drive.
\end{keyidea}

\begin{physicalbox}[Why oscillators dominate physics]
Near a stable equilibrium, a smooth potential is approximately quadratic:
\[
V(x)\approx V(x_*)+\frac12V''(x_*)(x-x_*)^2.
\]
Therefore small deviations obey an oscillator equation. Molecular vibrations, lattice modes, electrical resonances, optical cavities, and quantum field modes all inherit this local quadratic structure.
\end{physicalbox}
